\section{Introduction}
The repeated use of a measurement distinguishes perturbations that have
the same one-use size. Near a projection, a change of eigenspace can be
resolved on a scale of order $N^{-1}$, while an interior spectral
perturbation has scale $N^{-1/2}$. These scales meet at singular effects
and at repeated eigenvalues. We determine their joint geometry for the
entire body of ordered binary measurements in every fixed finite input
dimension, and derive the resulting description entropy.

Let
\[
 \mathfrak E_d=\{E=E^*:0\preceq E\preceq I_d\},\qquad
 \mathcal M_E(\rho)=\operatorname{tr}(E\rho)|1\rangle\langle1|
       +\operatorname{tr}((I_d-E)\rho)|0\rangle\langle0|.
\]
The output labels are ordered and the device has no residual quantum
output. For two such memoryless devices, $d_N$ is the supremum of the
unhalved trace distance between the final joint states of a common
experiment using at most $N$ calls. The experiment may retain an
arbitrary reference, entangle inputs across calls, use feedback, and
stop publicly. In particular $0\le d_N\le2$. The one-use identity
$d_1(\mathcal M_E,\mathcal M_F)=2\|E-F\|_{\rm op}$ is known
\cite[Lemmas III.7--III.8]{MB67}. We write $d_N^{\rm na}$ for the
restriction to nonadaptive experiments with the same reference access.

\subsection{The matrix comparison}
For $M=(E+F)/2$, $H=F-E$, and $V=M(I_d-M)$, define
\[
 Q_N(E,F)^2
   =N\langle H,(L_V+R_V+N^{-1}\operatorname{Id})^{-1}H
                      \rangle_{\rm HS},
\]
where $L_V(X)=VX$ and $R_V(X)=XV$. The inverse acts on the real
Hilbert space of Hermitian matrices. Our first main result is
Theorem~\ref{thm:matrixmetric76}:
\[
 \frac{\min\{1,Q_N(E,F)\}}{8192d}
 \le d_N^{\rm na}(\mathcal M_E,\mathcal M_F)
 \le d_N(\mathcal M_E,\mathcal M_F)
 \le\min\{2,8Q_N(E,F)\}.
\]
This comparison is uniform in $N$, in the effects, and in their
spectral multiplicities. It gives a horizon-independent comparison
between adaptive and nonadaptive discrimination at each fixed $d$;
it does not assert equality of their optimum values. The expression
$Q_N$ is an intrinsic comparison modulus, with no asserted triangle
inequality.

The upper bound uses a horizontal dilation of a binary measurement.
A regularized Sylvester equation splits an arbitrary matrix tangent
into a horizontal part and a remainder controlled by the hybrid
inequality. Operator concavity of $G(I_d-G)$ then integrates the
tangent estimate into the displayed midpoint formula on the closed
body. The lower bound uses one-dimensional Bernoulli tests and
two-dimensional compressions, whose full biased geometry is proved
in Section~\ref{sec:biasedmetric75}. This argument retains both
spectral endpoints and does not choose coordinates on a singular
eigenvector chart.

\subsection{Entropy and nested support scales}
Theorem~\ref{thm:matrixcover76} gives, for every fixed $d\ge1$,
\[
 \mathcal C_N(\mathfrak E_d,\delta)
 \asymp_d
 N^{d^2/2}[\log(N+2)]^{\lfloor d/2\rfloor}\delta^{-d^2},
 \qquad N\ge1,\quad0<\delta\le\delta_d.
\]
Here centres are arbitrary legal memoryless devices of the same
ordered binary interface; rational centres attain the same order.
The constant $\delta_d>0$ and both comparison constants depend only
on $d$. The corresponding fixed-length payload is
\[
 \frac{d^2}{2}\log_2N+\lfloor d/2\rfloor\log_2\log(N+2)
                +d^2\log_2(1/\delta)+O_d(1).
\]
The factor $[\log(N+2)]^{\lfloor d/2\rfloor}$ comes from nested
pairs of eigenvalues approaching zero and one. In the spectral
integral, the cumulative exponent of each ordered endpoint sector
is one half of the square of its signed prefix sum. Each balanced
prefix contributes a logarithm. The argument first bounds the
weighted volume of actual operational balls uniformly at every
boundary point; the spectral integral is used only after that
metric step has been established.

\subsection{A common learner and exact descriptions}
For an unknown effect in $\mathfrak E_2$, there is also a single
learning experiment with the uniform guarantee
\[
 \sup_{E\in\mathfrak E_2}
 \Pr_E\{d_N(\mathcal M_E,\mathcal M_{\widehat E})>\delta\}
 \le\eta.
\]
Theorem~\ref{thm:commonlearn76} proves that the optimum worst-case
training call count has order
$N\delta^{-2}\log(1/\eta)$ for small $\delta$ and
$0<\eta\le1/8$. Here $N$ is the future horizon in the loss.
Every entangled training block has length at most $N$.
Variance-sensitive product observations handle small contrast;
randomized GHZ experiments remove unknown bias and select their
block length from the data when the contrast is large. Fresh
endpoint observations complete the estimate, including at zero
eigenvalues. A scalar Bernoulli subfamily supplies the converse.

The rational code of Theorem~\ref{thm:biasedcodec75} converts this
estimate into one reusable word without additional device calls.
Its payload charges bias, contrast, and direction together. The
call bound, the payload length, and the classical cost of finding
the word are distinct resources. The arbitrary-dimensional
rational covering theorem asserts existence of optimal-order
codebooks; the implemented exact joint encoder treats qubits.

\subsection{Relation to prior work and organization}
Horizontal Kraus gauges and channel-extension metrology have
established antecedents \cite{FI76,DKG67,DM76,KGAD76}; integrating
channel distinguishability along a parameter path also predates
this work \cite{YF76,SD76}. Our matrix theorem provides the explicit
binary variance form, its uniform finite-horizon midpoint
comparison, and a matching nonadaptive lower bound. The spectral
volume calculation uses classical matrix integration
\cite{Dittmann76,SZGeom76}; its balanced-prefix accumulation gives
the logarithmic multiplicity above.

Finite-use measurement discrimination with entangled references
is developed in \cite{PPKKO72,KPP76,DBSA76}. The differing numbers
of outcomes and ancillary interfaces are compared in
Section~\ref{sec:matrixcomparison76}. The common learner uses
multiscale phase ideas \cite{HB76,KLY76} in a fully unknown biased
measurement experiment. Its future-horizon loss is distinct from
the one-use tomography guarantees of \cite{MB67,ZRK76}.

Sections~\ref{sec:matrixmetric76} and~\ref{sec:matrixentropy76}
prove the matrix comparison and entropy law. The common learning
theorem follows before the detailed qubit geometry and exact code.
The appendices retain the supporting instrument, preparation,
coherent, readout, and unbiased-boundary developments with their
complete proofs. The structural companion and the complete
research edition retain the stochastic-realization theory and
its arithmetic, profinite, and causal-process refinements.
